%MIT OpenCourseWare: https://ocw.mit.edu
%18.S190 Introduction to Metric Spaces, Independent Activities Period (IAP) 2023
%License: Creative Commons BY-NC-SA 
%For information about citing these materials or our Terms of Use, visit: https://ocw.mit.edu/terms.

\subsection*{Where We Go From Here}

In this section of the notes (the final one for this class!), we will discuss a bit of the history of metric spaces, and give a preview of how concepts learned here apply to future classes (i.e. 18.901, 18.102, etc.).

\textbf{History}: In the early 1900s, the usual approach to mathematics was far less abstract and axiomatic. Hence, at the time, \textit{various spaces} that mathematicians studied (such as function spaces as we have studied a bit of in this class) had different, ununified, notions of convergence. Each space has it's own notion of the word, which was studied in its own respect. There were some similarities between these notions, but there was no general understanding of the term. 

Then, in 1906, Fr\'echet introduced the idea of metric spaces in his \textit{Ph.D. Dissertation}. This allowed him (and many other mathematicians) to prove a result for a metric space, and have it be applicable to all other specific examples. This was the highlight of Lecture 2 of our class-- the General Theory of metric spaces.

\begin{remark}
Fr\'echet, however, did not coin the term "Metric Space"; the term was coined by Felix Hausdorff.
\end{remark}

In this class, we discussed \textit{normed spaces}, and saw how such spaces were in fact metric spaces under the metric induced by the norm. In this way, metric spaces are a generalization of normed spaces. 

\begin{question}
Is there a generalization of metric spaces?
\end{question}
Yes, there is: Topological spaces.

\textbf{18.901: Introduction to Topology}: Topological spaces were first defined by Hausdorff in 1914 in his book "Principles of Set Theory". His book increased the popularity of metric spaces as a mathematical tool. 

\begin{definition}[Topology]
    A \textbf{topology} on a set $X$ is a collection $\mathcal T$ of subsets of $X$ having the following properties:
\begin{enumerate}
    \item $\emptyset$ and $X$ are in $\mathcal T$.
    \item The union of the elements of any subcollection of $\mathcal T$ is in $\mathcal T$.
    \item The intersection of the elements of any finite subcollection of $\mathcal T$ is in $\mathcal T$.
\end{enumerate}
    A set $X$ for which a topology $\mathcal T$ has been specified is called a \textbf{topological space}.
\end{definition}

\begin{question}
Where has we seem properties like this before?
\end{question}

We saw these properties when we defined open sets! In fact, 
\begin{definition}[Open Set]
Let $(X, \mathcal T)$ be a topological space. Then, $U\subset X$ is called an \textbf{open set} if $U \in \mathcal T$. In other words, we define the set in $\mathcal T$ as open sets. Similarly, $V\subset X$ is a \textbf{closed set} if $X\setminus V \in \mathcal T$.
\end{definition}
Given a metric space $(X,d)$, there is a \textit{topology induced by the metric}. Intuitively, the topology is defined as the collection of unions of $\epsilon$-balls for all $\epsilon >0$. This intuitive definition follows from what we have discussed about metric spaces: an open set in a metric space is the union of (arbitrarily many) open balls.

\begin{remark}
To rigorously define the topology induced by the metric: $\epsilon$-balls form a basis for the topology on $X$.
\end{remark}

\begin{definition}[Metrizable]
    We say that a topological space $X$ is metrizable if there exists a metric $d$ on $X$ such that the topology on $X$ is the topology induced by $d$.
\end{definition}

\begin{remark}
    I briefly mentioned this on the second day, but again, not every topological space is metrizable!
\end{remark}

At first, the definition of a topology can feel too general to particularly seem useful, but then again the definition of metric spaces can feel the same way at first. Given a topological space $X$, we can define notions of neighborhoods, convergence, and continuity that align with what we have proven for metric spaces.
\begin{definition}
Let $(X,\mathcal T_X)$ and $(Y,\mathcal T_Y)$ be topological spaces. Then, 
\begin{enumerate}
    \item A \textbf{neighborhood} $U$ of a point $x\in X$ is an open set (i.e. $U \in \mathcal T_X$) such that $x\in U$.
    \item A sequence $\{x_n\}$ in $X$ \textbf{converges} to $x\in X$ if for every neighborhood $U$ of $x$, there is an $N$ such that $x_n \in U$ for all $n\geq N$.
    \item A function $f:X\to Y$ is said to be \textbf{continuous} if for each open set $V\in \mathcal T_Y$, $f^{-1}(V)\in \mathcal T_X$. 
\end{enumerate}
\end{definition}
This definition may feel very abstract, but as we have shown throughout this class, these definitions are related to our understanding of convergence and continuity in metric space.

One may ask, given that topological spaces are a generalization of metric spaces, why we study metric spaces. To this, I say: why do we study calculus before we study real analysis? In theory, we could prove calculus using real analysis without having actually taken 18.01 or 18.02. And yet, taking these pre-requisites give us an intuition for \textit{why} certain theorems should be true, and in some cases give us an intuition of how to approach proofs. Even more generally, calculus gives us an intuition for how derivatives of functions "should" look for nice enough functions. Even if we study weirder functions (or analogously weirder spaces that aren't metric spaces), having the intuition allows us to play with abstract objects.

This is why throughout this class I have tried to draw diagrams when possible to describe what a theorem is actually saying, or to visualize an example. This technique of drawing an entire "space" as a blob on a chalkboard or on paper is \textbf{very} useful when trying to approach a problem, which you can see more of if you decide to take a class like 18.901.

That being said, there is a \textit{lot} more we can know about metric spaces than we can about topological spaces, just in the same way we can know more about normed spaces than metric spaces. We have already seen this to a certain extent. One can show that given a normed space $(X, \lVert \cdot \rVert)$, and two convergent sequences $x_n\to x$ and $y_n \to y$, then $x_n+y_n \to x+y$. We cannot say the same about a general metric space, as we don't inherently have always have a notion of addition in a metric space. So metric spaces are still an interesting area to study even \textit{after} someone learns about topological spaces.  In fact, metric spaces are currently an active area of research, even though most research has ended for topological spaces.

So so far, we have talked about topological spaces which are more general than metric spaces. Is there a space that is more specific than a metric space that is useful to study? Yes, and we have already discussed it briefly: normed spaces.

\textbf{18.102: Introduction to Functional Analysis}: In Lecture 3 for this class, as a way to motivate the usefulness of compact sets, we defined normed spaces.

\begin{recall}
A vector space with a norm on it is defined as a \textbf{normed space}.
\end{recall}

And as you showed in PSET 2, a normed space is in fact a metric space, by defining $d(x,y) = \lVert x-y\rVert$: the metric induced by the norm. In fact, using our concept of convergence, Cauchy sequences, and open sets for metric spaces, we can define these ideas once again for normed spaces.
\begin{definition}
Let $(X,\lVert\cdot \rVert)$ be a normed space and $\{x_n\}$ be a sequence in $X$. Let $d$ be the metric induced by the norm. Then, 
\begin{enumerate}
    \item $x_n$ \textbf{converges} to $x$ if and only if for all $\epsilon >0$, there exists an $N$ such that for all $n\geq N$, 
    \[
    d(x_n,x) = \lVert x_n - x\rVert < \epsilon.
    \]
    \item $\{x_n\}$ is a \textbf{Cauchy sequence} if and only if for all $\epsilon >0$, there exists an $N$ such that for all $n,m \geq N$, 
    \[
    d(x_n, x_m) = \lVert x_n - x_m\rVert < \epsilon.
    \]
    \item A set $A$ is \textbf{open} in $X$ if for all $x\in A$, there exists an $\epsilon >0$ such that 
    \[
    B_\epsilon(x) = \{y\in X \mid d(x,y) = \lVert x-y\rVert < \epsilon\} \subset A.
    \]
\end{enumerate}
\end{definition}

We similarly have a definition of Cauchy completeness in a normed space (i.e. a space is Cauchy complete if every Cauchy sequence converges in the space). To be honest, when I first learned about Cauchy sequences real analysis, Cauchy sequences did not seem all that useful. So, I found it somewhat shocking that Cauchy sequences were extremely important in functional analysis.

\begin{definition}[Banach Space]
A \textbf{Banach space} is a normed space that is Cauchy complete with respect to the norm.
\end{definition}
Banach spaces are named after Stefan Banach who studied the spaces in 1920-1922.

\begin{remark}
Fun fact: The term "Banach space" was coined by Fr\'echet, and the term "Fr\'echet space" (which we did not cover in this class) was coined by Banach.
\end{remark}

\begin{example}
As we have shown (either in 18.100A/P or in this class), $\R^n$, $\CC^n$, and $C^0([a,b])$ are Banach spaces. In fact, one can show that the space 
\[
C_\infty(X) = \{f:X \to \CC \mid f \text{ continuous and bounded}\}
\]
is a Banach space with respect to the uniform norm on metric spaces.
\end{example}

\begin{definition}[Functional]
Let $(V,\lVert \cdot \rVert)$ be a normed space. A \textbf{functional} is a bounded linear map from $f: V\to \mathbb K$ where $\mathbb K = \R$ or $\CC$ depending on the context.
\end{definition}

This term is directly related to why 18.102 is called functional analysis! As it turns out, studying normed spaces is heavily related to studying functionals on those spaces. If you have taken linear algebra, this idea is similar to how studying a vector space is heavily related to studying the dual space of that vector space, but I digress.

\begin{example}
    One can show, which some work, that the set of functionals, $T$, are Cauchy complete under the operator norm, given by 
    \[
    \lVert  T\rVert_{\text{op}} = \sup{v\in V; \lVert v \rVert = 1} |Tv|.
    \]
    One can similarly define this operator norm for bounded linear maps \textit{between} normed spaces.
\end{example}

\begin{remark}
    18.102 is also heavily related to quantum mechanics, so if you are someone interested in physics/this concept, a class like functional analysis may be interesting to take.
\end{remark}

\textbf{18.101: Analysis and Manifolds}
I briefly want to mention how this class relates to (smooth) manifold theory. You can loosely think of a manifold as a smooth enough blob out in space that is locally flat (i.e. close up to the blob, it looks like Euclidean space). Some classic examples of manifolds include the sphere and open subsets of Euclidean space.

In the rigorous definition of a manifold, we suppose that the space is metrizable! We don't often \textit{directly} use the metric itself, but rather we use the fact that $\epsilon$-balls are open sets on the manifold. From this concept, you can develop the theory of/study smooth functions, integration, vector fields, and much much more on manifolds. Essentially, the idea is to develop a concept of multivariable calculus that is compatible to calculus on Euclidean space, but for manifolds!

\textbf{18.152: Introduction to Partial Differential Equations}: I can't say too much about this class as I haven't taken it myself, but just to bring it up again, the Banach fixed point theorem is very useful for proofs in this class regarding differential equations as we have seen.

On that note, I want to discuss one specific problem in partial differential equations that actually motivated the development of material in 18.100x classes: the Dirichlet problem. Consider some subset $\Omega \subset \R^2$, and picture this set as a metal plate. Suppose I heated the plate with a blowtorch for a certain amount of time. After a long period of time, the plate will reach thermal equilibrium. This is represented by the steady-state heat equation:
\[
\frac{\partial^2 u }{\partial x^2} + \frac{\partial^2 u }{\partial y^2} = 0
\]
where $u(x,y)$ is the temperature at point $(x,y).$ This differential operator is so important, we abbreviate it with $\Delta$, referred to as the Laplace operator or \textbf{Laplacian}. Suppose that I also knew the temperature on the boundary of this plate (denoted $\partial \Omega$). Let $u = f$ on $\partial \Omega$ (i.e. $f$ is the temperature on the boundary of the plate).

\begin{question}
Can we find the temperature distribution on $\Omega$? Is this temperature distribution unique?
\end{question}

 This question is known as the Dirichlet problem (in $\R^2$). To restate the question, we want to  solve the differential equation 
\[
\begin{cases}
\Delta u = 0 &\text{on } \Omega \\
u\bigr|_{\partial \Omega} = f
\end{cases}
\]
for some $u\in C^2(\R^2)$ and $f\in C^0(\R^2)$.

Early on, mathematicians studied this problem by studying the "Dirichlet energy" of $u$, given by the function 
\[
E(u) = \frac{1}{2}\int_\Omega
 |\nabla u|^2 \,\mathrm dA.
\]
They considered the set of all the $C^2(\Omega)$ with the given boundary condition and let $E_{\text{inf}}$ be the infimal energy of functions in this set. Dirichlet and others showed that if $E_{\text{inf}}$ is the energy of a function $u\in C^2$ with the given boundary condition, then $\Delta u = 0$. Thus, $u$ would be a solution to the Dirichlet problem! However, this raises the question:
\begin{question}[Question 1]
Does there exist a $C^2$ function $u$ with the given boundary condition, such that $E(u) = E_{\text{inf}}$?
\end{question}

If the answer is yes, then we have solved the problem! How might we approach finding such a function $u$? We could take a sequence of functions $u_n$ such that $E(u_n) \to E(u)$! However, this raises yet another question:
\begin{question}[Question 2]
Does $u_n$ converge in $C^2$ to a limit function $u$?
\end{question}

If this is true, then we can easily show that $E(u) = E_{\text{inf}}$ and then we have solve the problem! In fact, if $u_n$ converges in $C^1$ to $u$, then $E(u)=E_{\text{inf}}$. This raises the following question again:
\begin{question}[Question 3]
Consider the set of functions $u\in C^2$ with the given boundary data, with $E(u)$ at most $E_0>0$. Is this set of functions compact in $C^2$ or $C^1$? 
\end{question}
If \textit{this} is true, then we have sequential compactness, which means we can find a subsequence $u_{n_k}$ converging to a limit $u$ with $E(u) = E_{\text{inf}}$.

The answer to question 3, is no. This set of functions it \textit{not} compact in either $C^1$ or $C^2$. This ruins this entire, arguably very intuitive approach to the problem. Even though this idea didn't solve the problem though, it did highlight some key issues at play. In fact, the ideas of convergence and compactness in a metric space were developed partly to see if this approach to the Dirichlet problem works or not. It turns out, this approach doesn't work.

The answer to question 2 is also no. We can in fact build a sequence of $C^2$ functions $u_n$ with the given boundary condition such that $E(u_n) \to E_{\text{inf}}$, and yet $u_n$ fails to converge in $C^1$ (let alone $C^2)$. One can show that a function with zero boundary data and very small Dirichlet energy can \textit{still} have a large $C^1$ norm.

However, the answer to question 1 is \textit{yes}. Eventually, mathematicians were able to solve the Dirichlet problem, using techniques beyond the scope of this class. While we won't discuss the solution, this rich history highlights how mathematical concepts were/are developed to solve problems like this. This topic is sometimes talked about more in 18.102, and in some graduate level classes.

Thank you for taking this class with me this IAP. It has been a fun time developing the material and teaching the class. Please feel free to send me any feedback by emailing me or talking to me after class. Have a great spring semester!